\documentclass[10pt,a4paper]{amsart}
\usepackage{amsmath,amssymb,amsthm,graphicx}
% For commutative diagrams we use Xy-pic
\usepackage[all]{xy}

% We use 2cell for 2-commutative diagrams.
\xyoption{2cell}
% Theorem environments.
%
\theoremstyle{plain}
\newtheorem{theorem}[subsection]{Theorem}
\newtheorem{proposition}[subsection]{Proposition}
\newtheorem{lemma}[subsection]{Lemma}

\theoremstyle{definition}
\newtheorem{definition}[subsection]{Definition}
\newtheorem{example}[subsection]{Example}
\newtheorem{exercise}[subsection]{Exercise}
\newtheorem{situation}[subsection]{Situation}

\theoremstyle{remark}
\newtheorem{remark}[subsection]{Remark}
\newtheorem{remarks}[subsection]{Remarks}

\numberwithin{equation}{subsection}

% Macros
%
\def\lim{\mathop{\mathrm{lim}}\nolimits}
\def\colim{\mathop{\mathrm{colim}}\nolimits}
\def\Spec{\mathop{\mathrm{Spec}}}
\def\Hom{\mathop{\mathrm{Hom}}\nolimits}
\def\Ext{\mathop{\mathrm{Ext}}\nolimits}
\def\SheafHom{\mathop{\mathcal{H}\!\mathit{om}}\nolimits}
\def\SheafExt{\mathop{\mathcal{E}\!\mathit{xt}}\nolimits}
\def\Sch{\mathit{Sch}}
\def\Mor{\mathop{\mathrm{Mor}}\nolimits}
\def\Ob{\mathop{\mathrm{Ob}}\nolimits}
\def\Sh{\mathop{\mathit{Sh}}\nolimits}
\def\NL{\mathop{N\!L}\nolimits}
\def\CH{\mathop{\mathrm{CH}}\nolimits}
\def\proetale{{pro\text{-}\acute{e}tale}}
\def\etale{{\acute{e}tale}}
\def\QCoh{\mathit{QCoh}}
\def\Ker{\mathop{\mathrm{Ker}}}
\def\Im{\mathop{\mathrm{Im}}}
\def\Coker{\mathop{\mathrm{Coker}}}
\def\Coim{\mathop{\mathrm{Coim}}}

% Boxtimes
%
\DeclareMathSymbol{\boxtimes}{\mathbin}{AMSa}{"02}

%
% Macros for moduli stacks/spaces
%
\def\QCohstack{\mathcal{QC}\!\mathit{oh}}
\def\Cohstack{\mathcal{C}\!\mathit{oh}}
\def\Spacesstack{\mathcal{S}\!\mathit{paces}}
\def\Quotfunctor{\mathrm{Quot}}
\def\Hilbfunctor{\mathrm{Hilb}}
\def\Curvesstack{\mathcal{C}\!\mathit{urves}}
\def\Polarizedstack{\mathcal{P}\!\mathit{olarized}}
\def\Complexesstack{\mathcal{C}\!\mathit{omplexes}}
% \Pic is the operator that assigns to X its picard group, usage \Pic(X)
% \Picardstack_{X/B} denotes the Picard stack of X over B
% \Picardfunctor_{X/B} denotes the Picard functor of X over B
\def\Pic{\mathop{\mathrm{Pic}}\nolimits}
\def\Picardstack{\mathcal{P}\!\mathit{ic}}
\def\Picardfunctor{\mathrm{Pic}}
\def\Deformationcategory{\mathcal{D}\!\mathit{ef}}
\newenvironment{reference}{}{}
\newenvironment{slogan}{}{}
\newenvironment{history}{}{}
\usepackage{cleveref}
\crefname{originallemma}{original lemma}{original lemmas}
\Crefname{originallemma}{Original lemma}{Original lemmas}
\makeatletter
\newlabel{lemma-mod-injective}{{00ME}{source}{}{sourcepointer-1}{}}
\newlabel{lemma-mod-injective@cref}{{[originallemma][0][]00ME}{[1][0][]source}{}{}{}}
\makeatother
\crefname{originallemma}{original lemma}{original lemmas}
\Crefname{originallemma}{Original lemma}{Original lemmas}
\makeatletter
\newlabel{lemma-NAK}{{00DV}{source}{}{sourcepointer-2}{}}
\newlabel{lemma-NAK@cref}{{[originallemma][0][]00DV}{[1][0][]source}{}{}{}}
\makeatother
\crefname{originalproposition}{original proposition}{original propositions}
\Crefname{originalproposition}{Original proposition}{Original propositions}
\makeatletter
\newlabel{proposition-what-exact}{{00N1}{source}{}{sourcepointer-3}{}}
\newlabel{proposition-what-exact@cref}{{[originalproposition][0][]00N1}{[1][0][]source}{}{}{}}
\makeatother
\crefname{originallemma}{original lemma}{original lemmas}
\Crefname{originallemma}{Original lemma}{Original lemmas}
\makeatletter
\newlabel{lemma-open-regular-sequence}{{00RA}{source}{}{sourcepointer-4}{}}
\newlabel{lemma-open-regular-sequence@cref}{{[originallemma][0][]00RA}{[1][0][]source}{}{}{}}
\makeatother
\crefname{originallemma}{original lemma}{original lemmas}
\Crefname{originallemma}{Original lemma}{Original lemmas}
\makeatletter
\newlabel{lemma-limit-module-finite-presentation}{{00R1}{source}{}{sourcepointer-5}{}}
\newlabel{lemma-limit-module-finite-presentation@cref}{{[originallemma][0][]00R1}{[1][0][]source}{}{}{}}
\makeatother
\crefname{originallemma}{original lemma}{original lemmas}
\Crefname{originallemma}{Original lemma}{Original lemmas}
\makeatletter
\newlabel{lemma-limit-module-essentially-finite-presentation}{{00QX}{source}{}{sourcepointer-6}{}}
\newlabel{lemma-limit-module-essentially-finite-presentation@cref}{{[originallemma][0][]00QX}{[1][0][]source}{}{}{}}
\makeatother
\crefname{originallemma}{original lemma}{original lemmas}
\Crefname{originallemma}{Original lemma}{Original lemmas}
\makeatletter
\newlabel{lemma-colimit-eventually-flat}{{00R6}{source}{}{sourcepointer-7}{}}
\newlabel{lemma-colimit-eventually-flat@cref}{{[originallemma][0][]00R6}{[1][0][]source}{}{}{}}
\makeatother
\crefname{originallemma}{original lemma}{original lemmas}
\Crefname{originallemma}{Original lemma}{Original lemmas}
\makeatletter
\newlabel{lemma-flat-base-change}{{00HI}{source}{}{sourcepointer-8}{}}
\newlabel{lemma-flat-base-change@cref}{{[originallemma][0][]00HI}{[1][0][]source}{}{}{}}
\makeatother
\crefname{originallemma}{original lemma}{original lemmas}
\Crefname{originallemma}{Original lemma}{Original lemmas}
\makeatletter
\newlabel{lemma-resolution-by-finite-free}{{00LP}{source}{}{sourcepointer-9}{}}
\newlabel{lemma-resolution-by-finite-free@cref}{{[originallemma][0][]00LP}{[1][0][]source}{}{}{}}
\makeatother
\crefname{originallemma}{original lemma}{original lemmas}
\Crefname{originallemma}{Original lemma}{Original lemmas}
\makeatletter
\newlabel{lemma-flat-ses}{{00HM}{source}{}{sourcepointer-10}{}}
\newlabel{lemma-flat-ses@cref}{{[originallemma][0][]00HM}{[1][0][]source}{}{}{}}
\makeatother
\title{Openness of the flat locus of a finitely presented module over an arbitrary base}
\author{The Stacks Project Authors}
\date{stacks-project:algebra.tex:local-snapshot-20261011; GNU FDL 1.2 or later}
\begin{document}
\maketitle
\noindent\textit{Editorial note.} Public1510 assumes locally finitely many weakly associated base points and a finite-type sheaf; it uses the usual finitely-presented openness theorem as input. This task proves the latter for arbitrary base rings. Neither hypothesis set is merely the same renamed assertion.\par\medskip
\noindent\textit{Editorial note.} The source describes the localized polynomial fibre global dimension as n; the needed bound is at most n, which suffices for the nth syzygy. One may add dummy polynomial variables to avoid the harmless n=0 indexing convention.\par\medskip
\noindent\textit{Allowed standard prerequisites (editorial).}\begin{itemize}
\item Explicit prerequisite boundary. Finite-presentation approximation and descent of flatness, Tor exactness, polynomial-ring global dimension, local flat/free criterion and finite free resolutions. The local fibrewise-exactness and final cokernel arguments are retained as author proof, not removed.
\end{itemize}\medskip
\ifdefined\hypertarget\hypertarget{sourcepointer-1}{}\fi
\noindent\textit{Source reference.} Original lemma 00ME: focused-seeds/AGP/sources/stacks-snapshot-20261011/tags.txt, source lines 820--820. Stable Stacks tag for an expressly allowed foundational or notational input; the counted construction remains the two concrete core obligations. The retained local core is not replaced by this pointer. See standard\_prerequisites.\par\medskip
\ifdefined\hypertarget\hypertarget{sourcepointer-2}{}\fi
\noindent\textit{Source reference.} Original lemma 00DV: focused-seeds/AGP/sources/stacks-snapshot-20261011/tags.txt, source lines 512--512. Stable Stacks tag for an expressly allowed foundational or notational input; the counted construction remains the two concrete core obligations. The retained local core is not replaced by this pointer. See standard\_prerequisites.\par\medskip
\ifdefined\hypertarget\hypertarget{sourcepointer-3}{}\fi
\noindent\textit{Source reference.} Original proposition 00N1: focused-seeds/AGP/sources/stacks-snapshot-20261011/tags.txt, source lines 843--843. Stable Stacks tag for an expressly allowed foundational or notational input; the counted construction remains the two concrete core obligations. The retained local core is not replaced by this pointer. See standard\_prerequisites.\par\medskip
\ifdefined\hypertarget\hypertarget{sourcepointer-4}{}\fi
\noindent\textit{Source reference.} Original lemma 00RA: focused-seeds/AGP/sources/stacks-snapshot-20261011/tags.txt, source lines 996--996. Stable Stacks tag for an expressly allowed foundational or notational input; the counted construction remains the two concrete core obligations. The retained local core is not replaced by this pointer. See standard\_prerequisites.\par\medskip
\ifdefined\hypertarget\hypertarget{sourcepointer-5}{}\fi
\noindent\textit{Source reference.} Original lemma 00R1: focused-seeds/AGP/sources/stacks-snapshot-20261011/tags.txt, source lines 987--987. Stable Stacks tag for an expressly allowed foundational or notational input; the counted construction remains the two concrete core obligations. The retained local core is not replaced by this pointer. See standard\_prerequisites.\par\medskip
\ifdefined\hypertarget\hypertarget{sourcepointer-6}{}\fi
\noindent\textit{Source reference.} Original lemma 00QX: focused-seeds/AGP/sources/stacks-snapshot-20261011/tags.txt, source lines 983--983. Stable Stacks tag for an expressly allowed foundational or notational input; the counted construction remains the two concrete core obligations. The retained local core is not replaced by this pointer. See standard\_prerequisites.\par\medskip
\ifdefined\hypertarget\hypertarget{sourcepointer-7}{}\fi
\noindent\textit{Source reference.} Original lemma 00R6: focused-seeds/AGP/sources/stacks-snapshot-20261011/tags.txt, source lines 992--992. Stable Stacks tag for an expressly allowed foundational or notational input; the counted construction remains the two concrete core obligations. The retained local core is not replaced by this pointer. See standard\_prerequisites.\par\medskip
\ifdefined\hypertarget\hypertarget{sourcepointer-8}{}\fi
\noindent\textit{Source reference.} Original lemma 00HI: focused-seeds/AGP/sources/stacks-snapshot-20261011/tags.txt, source lines 644--644. Stable Stacks tag for an expressly allowed foundational or notational input; the counted construction remains the two concrete core obligations. The retained local core is not replaced by this pointer. See standard\_prerequisites.\par\medskip
\ifdefined\hypertarget\hypertarget{sourcepointer-9}{}\fi
\noindent\textit{Source reference.} Original lemma 00LP: focused-seeds/AGP/sources/stacks-snapshot-20261011/tags.txt, source lines 795--795. Stable Stacks tag for an expressly allowed foundational or notational input; the counted construction remains the two concrete core obligations. The retained local core is not replaced by this pointer. See standard\_prerequisites.\par\medskip
\ifdefined\hypertarget\hypertarget{sourcepointer-10}{}\fi
\noindent\textit{Source reference.} Original lemma 00HM: focused-seeds/AGP/sources/stacks-snapshot-20261011/tags.txt, source lines 648--648. Stable Stacks tag for an expressly allowed foundational or notational input; the counted construction remains the two concrete core obligations. The retained local core is not replaced by this pointer. See standard\_prerequisites.\par\medskip
\begin{theorem}
\label{theorem-openness-flatness}
Let $R$ be a ring. Let $R \to S$ be a ring map of finite
presentation. Let $M$ be a finitely presented $S$-module.
The set
$$
\{ \mathfrak q \in \Spec(S) \mid
M_{\mathfrak q}\text{ is flat over }R\}
$$
is open in $\Spec(S)$.
\end{theorem}

\begin{lemma}
\label{lemma-free-fibre-flat-free}
Let $R \to S$ be a local homomorphism of Noetherian
local rings. Let $\mathfrak m$ be the maximal
ideal of $R$. Let $M$ be a nonzero finite $S$-module.
Suppose that (a) $M/\mathfrak mM$
is a free $S/\mathfrak mS$-module, and (b) $M$ is flat over $R$.
Then $M$ is free and $S$ is flat over $R$.
\end{lemma}

\begin{proof}
Let $\overline{x}_1, \ldots, \overline{x}_n$ be a basis
for the free module $M/\mathfrak mM$. Choose
$x_1, \ldots, x_n \in M$ with $x_i$ mapping to $\overline{x}_i$. Let
$u : S^{\oplus n} \to M$ be the map which maps the $i$th
standard basis vector to $x_i$. By Lemma \ref{lemma-mod-injective}
we see that $u$ is injective. On the other hand, by
Nakayama's Lemma \ref{lemma-NAK} the map is surjective. The
lemma follows.
\end{proof}

\begin{lemma}
\label{lemma-complex-exact-mod}
Let $R \to S$ be a local homomorphism of local Noetherian
rings. Let $\mathfrak m$ be the maximal ideal of $R$.
Let $0 \to F_e \to F_{e-1} \to \ldots \to F_0$
be a finite complex of finite $S$-modules. Assume that
each $F_i$ is $R$-flat, and that the complex
$0 \to F_e/\mathfrak m F_e \to F_{e-1}/\mathfrak m F_{e-1}
\to \ldots \to F_0 / \mathfrak m F_0$ is exact.
Then $0 \to F_e \to F_{e-1} \to \ldots \to F_0$
is exact, and moreover the module
$\Coker(F_1 \to F_0)$ is $R$-flat.
\end{lemma}

\begin{proof}
By induction on $e$. If $e = 1$, then this is exactly
Lemma \ref{lemma-mod-injective}. If $e > 1$, we see
by Lemma \ref{lemma-mod-injective} that $F_e \to F_{e-1}$
is injective and that $C = \Coker(F_e \to F_{e-1})$
is a finite $S$-module flat over $R$. Hence we can
apply the induction hypothesis to the complex
$0 \to C \to F_{e-2} \to \ldots \to F_0$.
We deduce that $C \to F_{e-2}$ is injective
and the exactness of the complex follows, as well
as the flatness of the cokernel of $F_1 \to F_0$.
\end{proof}

\begin{lemma}
\label{lemma-exact-on-fibres-open}
Let $R \to S$ be a ring map. Consider a finite homological complex of
finite free $S$-modules:
$$
F_{\bullet} :
0
\to
S^{n_e}
\xrightarrow{\varphi_e}
S^{n_{e-1}}
\xrightarrow{\varphi_{e-1}}
\ldots
\xrightarrow{\varphi_{i + 1}}
S^{n_i}
\xrightarrow{\varphi_i}
S^{n_{i-1}}
\xrightarrow{\varphi_{i-1}}
\ldots
\xrightarrow{\varphi_1}
S^{n_0}
$$
For every prime $\mathfrak q$ of $S$ consider the
complex $\overline{F}_{\bullet, \mathfrak q} =
F_{\bullet, \mathfrak q} \otimes_R \kappa(\mathfrak p)$
where $\mathfrak p$ is inverse image of $\mathfrak q$ in $R$.
Assume $R$ is Noetherian and there exists an integer $d$ such
that $R \to S$ is finite type, flat
with fibres $S \otimes_R \kappa(\mathfrak p)$
Cohen-Macaulay of dimension $d$.
The set
$$
\{\mathfrak q \in \Spec(S) \mid
\overline{F}_{\bullet, \mathfrak q}\text{ is exact}\}
$$
is open in $\Spec(S)$.
\end{lemma}

\begin{proof}
Let $\mathfrak q$ be an element of the set defined in the lemma.
We are going to use Proposition \ref{proposition-what-exact}
to show there exists a $g \in S$, $g \not \in \mathfrak q$
such that $D(g)$ is contained in the set defined in the lemma.
In other words, we are going to show that after replacing $S$
by $S_g$, the set of the lemma is all of $\Spec(S)$.
Thus during the proof we will, finitely often, replace
$S$ by such a localization.
Recall that Proposition \ref{proposition-what-exact}
characterizes exactness of complexes
in terms of ranks of the maps $\varphi_i$ and the ideals
$I(\varphi_i)$, in case the ring is local. We first address
the rank condition. Set
$r_i = n_i - n_{i + 1} + \ldots + (-1)^{e - i} n_e$.
Note that $r_i + r_{i + 1} = n_i$ and note that
$r_i$ is the expected rank of $\varphi_i$ (in the
exact case).

\medskip\noindent
By Lemma \ref{lemma-complex-exact-mod} we see that if
$\overline{F}_{\bullet, \mathfrak q}$ is exact, then
the localization $F_{\bullet, \mathfrak q}$ is exact.
In particular the complex $F_\bullet$ becomes
exact after localizing by an element
$g \in S$, $g \not \in \mathfrak q$. In this case
Proposition \ref{proposition-what-exact} applied
to all localizations of $S$ at prime ideals
implies that all $(r_i + 1) \times (r_i + 1)$-minors
of $\varphi_i$ are zero. Thus we see that the rank
of $\varphi_i$ is at most $r_i$.

\medskip\noindent
Let $I_i \subset S$ denote the ideal generated
by the $r_i \times r_i$-minors of the matrix
of $\varphi_i$. By Proposition \ref{proposition-what-exact}
the complex $\overline{F}_{\bullet, \mathfrak q}$ is exact
if and only if for every $1 \leq i \leq e$ we have
either $(I_i)_{\mathfrak q} = S_{\mathfrak q}$ or
$(I_i)_{\mathfrak q}$ contains a $S_{\mathfrak q}/\mathfrak p
S_{\mathfrak q}$-regular sequence of length $i$.
Namely, by our choice of $r_i$ above and by the
bound on the ranks of the $\varphi_i$ this is the
only way the conditions of Proposition \ref{proposition-what-exact}
can be satisfied.

\medskip\noindent
If $(I_i)_{\mathfrak q} = S_{\mathfrak q}$, then after localizing $S$ at
some element $g \not\in \mathfrak q$ we may assume that
$I_i = S$. Clearly, this is an open condition.

\medskip\noindent
If $(I_i)_{\mathfrak q} \not = S_{\mathfrak q}$, then we have
a sequence $f_1, \ldots, f_i \in (I_i)_{\mathfrak q}$ which
form a regular sequence in $S_{\mathfrak q}/\mathfrak pS_{\mathfrak q}$.
Note that for any prime $\mathfrak q' \subset S$ such that
$(f_1, \ldots, f_i) \not \subset \mathfrak q'$ we have
$(I_i)_{\mathfrak q'} = S_{\mathfrak q'}$.
Thus the result follows from Lemma \ref{lemma-open-regular-sequence}.
\end{proof}

\begin{proof}
Let $\mathfrak q \in \Spec(S)$ be a prime.
Let $\mathfrak p \subset R$ be the inverse image of $\mathfrak q$ in $R$.
Note that $M_{\mathfrak q}$ is flat over $R$ if and only if
it is flat over $R_{\mathfrak p}$.
Let us assume that $M_{\mathfrak q}$ is flat over $R$.
We claim that there exists a $g \in S$, $g \not \in \mathfrak q$
such that $M_g$ is flat over $R$.

\medskip\noindent
We first reduce to the case where $R$ and $S$ are
of finite type over $\mathbf{Z}$.
Choose a directed set $\Lambda$ and
a system $(R_\lambda \to S_\lambda, M_\lambda)$
as in Lemma \ref{lemma-limit-module-finite-presentation}.
Set $\mathfrak p_\lambda$ equal to the inverse image of
$\mathfrak p$ in $R_\lambda$.
Set $\mathfrak q_\lambda$ equal to the inverse image of
$\mathfrak q$ in $S_\lambda$.
Then the system
$$
((R_\lambda)_{\mathfrak p_\lambda},
(S_\lambda)_{\mathfrak q_\lambda},
(M_\lambda)_{\mathfrak q_{\lambda}})
$$
is a system as in
Lemma \ref{lemma-limit-module-essentially-finite-presentation}.
Hence by Lemma \ref{lemma-colimit-eventually-flat}
we see that for some $\lambda$ the module
$M_\lambda$ is flat over $R_\lambda$ at the prime
$\mathfrak q_{\lambda}$. Suppose we
can prove our claim for the system
$(R_\lambda \to S_\lambda, M_\lambda, \mathfrak q_{\lambda})$.
In other words, suppose that we can find a $g \in S_\lambda$,
$g \not\in \mathfrak q_\lambda$ such that $(M_\lambda)_g$
is flat over $R_\lambda$. By Lemma \ref{lemma-limit-module-finite-presentation}
we have $M = M_\lambda \otimes_{R_\lambda} R$ and hence
also $M_g = (M_\lambda)_g \otimes_{R_\lambda} R$. Thus by
Lemma \ref{lemma-flat-base-change} we deduce the claim
for the system $(R \to S, M, \mathfrak q)$.

\medskip\noindent
At this point we may assume that $R$ and $S$ are of finite type
over $\mathbf{Z}$. We may write $S$ as a quotient of a
polynomial ring $R[x_1, \ldots, x_n]$. Of course, we may replace
$S$ by $R[x_1, \ldots, x_n]$ and assume that $S$ is a polynomial
ring over $R$. In particular we see that $R \to S$ is flat
and all fibres rings $S \otimes_R \kappa(\mathfrak p)$
have global dimension $n$.

\medskip\noindent
Choose a resolution $F_\bullet$ of $M$ over $S$ with each
$F_i$ finite free, see Lemma \ref{lemma-resolution-by-finite-free}.
Let $K_n = \Ker(F_{n-1} \to F_{n-2})$. Note that
$(K_n)_{\mathfrak q}$ is flat over $R$, since each $F_i$
is flat over $R$ and by assumption on $M$, see Lemma
\ref{lemma-flat-ses}. In addition, the sequence
$$
0 \to
K_n/\mathfrak p K_n \to
F_{n-1}/ \mathfrak p F_{n-1} \to
\ldots \to
F_0 / \mathfrak p F_0 \to
M/\mathfrak p M \to
0
$$
is exact upon localizing at $\mathfrak q$, because of vanishing
of $\text{Tor}_i^{R_\mathfrak p}(\kappa(\mathfrak p), M_{\mathfrak q})$.
Since the global dimension of $S_\mathfrak q/\mathfrak p S_{\mathfrak q}$
is $n$ we conclude that $K_n / \mathfrak p K_n$ localized
at $\mathfrak q$ is a finite free module over
$S_\mathfrak q/\mathfrak p S_{\mathfrak q}$. By
Lemma \ref{lemma-free-fibre-flat-free} $(K_n)_{\mathfrak q}$
is free over $S_{\mathfrak q}$. In particular, there exists a
$g \in S$, $g \not \in \mathfrak q$ such that $(K_n)_g$
is finite free over $S_g$.

\medskip\noindent
By Lemma \ref{lemma-exact-on-fibres-open}
there exists a further localization $S_g$ such that
the complex
$$
0 \to K_n \to F_{n-1} \to \ldots \to F_0
$$
is exact on {\it all fibres} of $R \to S$. By
Lemma \ref{lemma-complex-exact-mod}
this implies that the cokernel of $F_1 \to F_0$ is
flat. This proves the theorem in the Noetherian case.
\end{proof}

\bibliographystyle{plain}
\bibliography{my}
\end{document}
